\subsection{Structural Result Cluster: Three Types of Commutative Diagrams, Two Types of Indices, and Unique Continuous Transversal}\label{subsec:discussion-gc-three-diagrams-two-indices-unique-transversal}
This subsection further strictifies the ``fixed slice image recycler'' phenomenon from the previous subsection into a set of fully verifiable algebraic/operator invariants:
three types of strictly commutative diagrams compress visible information to a unified audit interface, two types of integer/dimensional indices characterize the discrete residuals allowed by the regime, and if one wishes to preserve continuous degrees of freedom outside the slice, the unique continuous transversal must factorize through \(\log|u|\) (Corollary \ref{cor:xi-unique-continuous-transverse-register}).
All statements below are anchored by previously numbered conclusions, serving only as structural consolidation and introducing no new proof premises.

\paragraph{Three types of commutative diagrams (compression mechanisms)}
The ``recycling'' in this paper's context can be precisely decomposed into three strictly commutative compression chains:
\begin{enumerate}
  \item \textbf{Witt--Euler commutative diagram (conjugacy forgetting):}
  By Theorem \ref{thm:discussion-witt-euler-gc-functor} (see also Remark \ref{rem:zetaK-commuting-diagram}),
  \[
  \mathrm{ZetaDet}
  ~=~
  \mathrm{Euler}\circ\mathrm{Mob}\circ\mathrm{TrSeq},
  \]
  i.e., any endomorphism within the regime is read out only through trace sequences/primitive sequences/Euler products and these readings are strictly equivalent.
  \item \textbf{Completed Chebyshev--Adams compression (power semigroup recycled as Adams operations on \(\ZZ[S]\)):}
  The power substitution \(u\mapsto u^d\) in completed trace coordinates \(S=r+r^{-1}\) is strictly compressed to \(S\mapsto C_d(S)\) (Appendix \S\ref{subsec:unit-disk-witt-chebyshev}, see also Corollary \ref{cor:chebyshev-witt-completed}), thereby placing all angle-doubling/completion audits within the same integer-coefficient closure.
  \item \textbf{Unitary slice crossing commutative diagram (zero event \(\Leftrightarrow\) crossing \(1\)):}
  Self-dual completion strictly compresses a critical line zero event to a spectral crossing event at the single point \(1\) on the unitary slice (Theorem \ref{thm:discussion-critical-line-single-point-crossing}).
\end{enumerate}

\begin{theorem}[Chebyshev--Witt equivariance: after completion, M\"obius/Witt inversion and Adams operations strictly commute]\label{thm:discussion-chebyshev-witt-equivariance}
Take the half-angle variable \(r\) and let \(u=r^2\), \(S=r+r^{-1}\in[-2,2]\).
For each \(d\ge 1\), define the Chebyshev--Adams polynomial
\[
C_d(S):=r^d+r^{-d}\in\ZZ[S].
\]
Let \(\{a_n(u)\}_{n\ge 1}\) be the trace (ghost) polynomial sequence and \(\{p_n(u)\}_{n\ge 1}\) its primitive sequence (Witt/M\"obius inversion, Proposition \ref{prop:zetaK-mobius-primitive}).
In completed coordinates, define the completed ghost and completed primitive:
\[
\widehat a_n(S):=r^{-n}a_n(r^2)\in\ZZ[S],\qquad
\widehat p_n(S):=r^{-n}p_n(r^2)\in\ZZ[S].
\]
Then for any \(n\ge 1\), the following strict identities hold:
\[
\boxed{\
\widehat p_n(S)=\frac1n\sum_{d\mid n}\mu(d)\,\widehat a_{n/d}\!\bigl(C_d(S)\bigr),
\qquad
\widehat a_n(S)=\sum_{d\mid n}d\,\widehat p_d\!\bigl(C_{n/d}(S)\bigr)\ }.
\]
Furthermore, the completed power substitution satisfies the strict ``angle-doubling law'':
\[
\boxed{\ r^{-dn}\,p_n(r^{2d})=\widehat p_n\!\bigl(C_d(S)\bigr)\ }.
\]
\end{theorem}

\begin{proof}
See Corollary \ref{cor:chebyshev-witt-completed} and Corollary \ref{cor:primitive-completion-hatp}; the descent of \(\widehat a_n,\widehat p_n\in\ZZ[S]\) follows from the inversion-invariant subring identity (Remark \ref{rem:completion-invariant-ring} or Proposition \ref{prop:completion-subring-unique-angle}). This completes the proof.
\end{proof}

\begin{corollary}[Conceptual compression proposition: primitive extraction is a \texorpdfstring{$\lambda$}{lambda}-equivariant structure on \texorpdfstring{$\ZZ[S]$}{Z[S]}]\label{cor:discussion-lambda-equivariant-primitive}
On \(\ZZ[S]\), define the Adams operation family \(\psi^d:\ZZ[S]\to\ZZ[S]\) (\(d\ge 1\)):
\[
\boxed{\ (\psi^d f)(S):=f\!\bigl(C_d(S)\bigr)\ }.
\]
Then the completed M\"obius/Witt inversion is strictly equivariant on \((\ZZ[S],\psi^d)\): \(\{\widehat a_n\}\) is the ghost data, \(\{\widehat p_n\}\) is the Big-Witt primitive block, and ``power semigroup recycling'' constitutes the auditable primitive coordinates on the same \(\lambda\)-structure.
\end{corollary}

\begin{proof}
Directly from the two reciprocal identities in Theorem \ref{thm:discussion-chebyshev-witt-equivariance}: all power substitutions enter only in the form \(S\mapsto C_d(S)\), which is the action of \(\psi^d\). This completes the proof.
\end{proof}

\begin{corollary}[Algebraic closure of the audit regime: angle-doubling/completion certificates on the unitary-slice can be fully realized in \texorpdfstring{$\ZZ[x]$}{Z[x]}]\label{cor:discussion-audit-closure-Zx}
Taking the Cayley coordinate \(r=\iota(x)=\frac{1+\mathrm{i}x}{1-\mathrm{i}x}\) (\(x\in\RR\)), then
\[
S(x)=r+r^{-1}=\frac{2(1-x^2)}{1+x^2},
\]
and for each \(d\ge 1\) there exists \(\tau_d(x)=B_d(x)/A_d(x)\in\QQ(x)\) (where \(A_d,B_d\in\ZZ[x]\) are generated by pure multiply-add recursion) such that
\[
C_d\!\bigl(S(x)\bigr)=S\!\bigl(\tau_d(x)\bigr).
\]
Therefore, any unitary-slice certificate chain depending only on \(S\mapsto C_d(S)\) (including iterated substitution and combination of \(\widehat a_n,\widehat p_n\)) can be fully realized within the multiply-add operations of \(\ZZ[x]\).
\end{corollary}

\begin{proof}
See Proposition \ref{prop:chebyshev-adams-fully-rational}. This completes the proof.
\end{proof}

\begin{theorem}[Completed Dwork congruence (Chebyshev form)]\label{thm:discussion-completed-dwork-chebyshev}
Let \(p\) be a prime and \(k\ge 1\). In the completed notation of Theorem \ref{thm:discussion-chebyshev-witt-equivariance}, the completed ghost satisfies the congruence
\[
\boxed{\ \widehat a_{p^k}(S)\equiv \widehat a_{p^{k-1}}\!\bigl(C_p(S)\bigr)\pmod{p^k}\qquad\text{in }\ZZ[S].\ }
\]
\end{theorem}

\begin{proof}
By Corollary \ref{cor:witt-dwork-congruence}, in \(\ZZ[u]\) we have
\(
a_{p^k}(u)\equiv a_{p^{k-1}}(u^p)\ (\mathrm{mod}\ p^k)
\).
Substituting \(u=r^2\) gives
\(
a_{p^k}(r^2)\equiv a_{p^{k-1}}(r^{2p})\ (\mathrm{mod}\ p^k)
\) in \(\ZZ[r,r^{-1}]\).
Multiplying both sides by \(r^{-p^k}\) and noting
\[
r^{-p^k}a_{p^{k-1}}(r^{2p})
=(r^p)^{-p^{k-1}}a_{p^{k-1}}\!\bigl((r^p)^2\bigr)
=\widehat a_{p^{k-1}}\!\bigl(S(r^p)\bigr)
=\widehat a_{p^{k-1}}\!\bigl(C_p(S)\bigr),
\]
we obtain the stated congruence; since \(\widehat a_{\bullet}\in\ZZ[S]\) (Theorem \ref{thm:discussion-chebyshev-witt-equivariance}), the congruence can be viewed as holding in \(\ZZ[S]\). This completes the proof.
\end{proof}

\begin{corollary}[Recycling residuals: mod \texorpdfstring{$p^k$}{p^k} guardrails and \texorpdfstring{$\mathbb{F}_p$}{F\_p} stable fingerprints]\label{cor:discussion-recycling-residuals}
Under the regime allowing only fixed-slice visible quantities, beyond the implementation details recycled by the commutative diagrams, two types of fully internal, offline-verifiable residual fingerprints remain:
\begin{enumerate}
  \item \textbf{\(\ZZ/p^k\)-congruence-type residuals (Dwork/necklace guardrails):}
  The completed Dwork congruence (Theorem \ref{thm:discussion-completed-dwork-chebyshev}) and its coefficient sparsification/iterated descent (Proposition \ref{prop:witt-pk-sparsification} and Proposition \ref{prop:witt-frobenius-iterated-descent}) provide level-by-level auditable \(\bmod p^k\) constraints.
  \item \textbf{\(\mathbb{F}_p\)-stable fingerprints:}
  The \(p\)-adic filtration of the \(\theta\)-derivative moments \(A^{(r)}_{p^k}\) produces stable leading-digit residuals \(\alpha_{p,r}\in\mathbb{F}_p\) independent of the level \(k\) (Corollary \ref{cor:witt-theta-stable-residue}); at \(p=2\), this further yields permanently stable binary fingerprints \(\varepsilon_r\in\{0,1\}\) (Corollary \ref{cor:witt-theta-binary-fingerprint}).
\end{enumerate}
Therefore, ``recycling residuals'' strictly fall into finite algebraic fingerprints: congruence guardrails \(\bmod p^k\) and stable fingerprints \(\bmod p\).
\end{corollary}

\paragraph{Two types of indices (discrete anomalies allowed by the regime)}
If ``off-slice degrees of freedom'' are understood as auditable differences that must be externalized as registers, then this paper yields two types of minimal indices:
\begin{enumerate}
  \item \textbf{Integer index \(\kappa\) of pointwise defects:}
  Under boundary unitarization, absence of singular inner factor, and finite defect caliber,
  \(\kappa=\mathrm{Wind}(F)=\deg(B_F)\)
  (Proposition \ref{prop:winding-equals-blaschke-degree}), and is unified equivalently with the Toeplitz negative square/Pontryagin/Tate--Maslov barrier (Theorem \ref{thm:xi-tate-maslov-pontryagin-unification}).
  \item \textbf{Dimensional \(\mathsf{NULL}\) index \(d_\ast\):}
  Under unified minimality and protocol closure, mode axis completeness holds if and only if the Hankel rank is uniformly bounded, i.e., there exists a uniform dimension upper bound \(d_\ast\) such that the minimal realization at all levels/phases does not exceed this dimension (Theorem \ref{thm:xi-completeness-iff-hankel-bounded}).
\end{enumerate}

\begin{theorem}[Unified index for crossing--winding--Blaschke degree: pointwise anomalies allowed by the regime can only appear in integer form]\label{thm:discussion-crossing-winding-blaschke-index}
Under self-dual completion and unitary-slice reduction, a critical line zero event is equivalent to the crossing event \(1\in\mathrm{Spec}(U_L(\theta))\) on the unitary slice (Theorem \ref{thm:discussion-critical-line-single-point-crossing}).
If furthermore one is in the boundary unitarization regime with no singular inner factor \(S_F\equiv 1\) and finite defect caliber, then there exists an integer index
\[
\boxed{\ \kappa=\mathrm{Wind}(F)=\deg(B_F)\ }
\]
and this \(\kappa\) completely characterizes the pointwise off-slice defects allowed within the regime (unification of Blaschke pointwise defect/Toeplitz failure mode/Tate--Maslov barrier, Theorem \ref{thm:xi-tate-maslov-pontryagin-unification}).
In particular, \(\deg(B_F)=0\) is the strict necessary condition for ``pointwise defect invisibilization (not externalizing an orthogonal residual axis).''
\end{theorem}

\begin{proof}
The first statement is Theorem \ref{thm:discussion-critical-line-single-point-crossing}.
Under the finite Blaschke and \(S_F\equiv 1\) caliber, Proposition \ref{prop:winding-equals-blaschke-degree} gives \(\mathrm{Wind}(F)=\deg(B_F)\);
and the pointwise defect as a regime invariant is unified into a single integer \(\kappa\) by Theorem \ref{thm:xi-tate-maslov-pontryagin-unification}.
The last statement follows immediately from \(\deg(B_F)=0\iff B_F\equiv 1\) (Remark \ref{rem:blaschke-defect-finite}). This completes the proof.
\end{proof}

\begin{theorem}[Unique continuous transversal direction: if continuous off-slice degrees of freedom are to be preserved, the only viable register must factorize through \texorpdfstring{$\log|u|$}{log|u|}]\label{thm:discussion-unique-continuous-transversal}
If one allows externalizing a continuous record axis beyond the pure-phase regime to carry radial degrees of freedom, then the continuous transversal direction is unique up to factorization: any continuous dimension extension that ``reduces \(\mathsf{NULL}\)'' must be equivalent to a monotone reparametrization of the externalized radial register \(\varrho(u)=\log|u|\); no second continuous transversal register parallel to the phase axis exists.
\end{theorem}

\begin{proof}
See Corollary \ref{cor:xi-unique-continuous-transverse-register} (in conjunction with the ``no additional continuous hair'' principle of Theorem \ref{thm:xi-pw-no-continuous-hair}). This completes the proof.
\end{proof}

\begin{corollary}[Discrete residuals and precision potential ledger at the endpoint: unique continuous channel \texorpdfstring{$\perp$}{⊥} finite discrete channels]\label{cor:discussion-endpoint-discrete-vs-continuous}
In the neighborhood of the endpoint \(u=-1\), the power semigroup action decomposes in completed coordinates into two mutually decoupled payloads:
\begin{enumerate}
  \item \textbf{Discrete endpoint residual:} \(C_m(0)\in\{0,\pm 2\}\) is determined solely by \(m\bmod 4\) (Proposition \ref{prop:half-angle-z4-residue});
  \item \textbf{Continuous precision ledger:} The endpoint precision potential \(y(u)=-\log|\mathfrak{w}(u)|\) satisfies the additive law \(y(\psi_m(u))=y(u)+\log m\) under dilation (Proposition \ref{prop:endpoint-orthogonal-decomposition-z4-Rplus}, see also \S\ref{subsec:unit-circle-scale-linearization}).
\end{enumerate}
Therefore, within the pure-phase closure that forbids a radial continuous register, the continuous information channel outside the slice is structurally sealed (can only trigger \(\mathsf{NULL}\)); the nontrivial residuals still allowed by the regime can only fall on countable/discrete channels (e.g., the integer index \(\kappa\) and the finite endpoint residual).
\end{corollary}

\begin{theorem}[Verifiable trigger condition for ``dimensional \texorpdfstring{$\mathsf{NULL}$}{NULL}'': Hankel rank uniform boundedness \texorpdfstring{$\Longleftrightarrow$}{<->} mode axis completeness]\label{thm:discussion-dimensional-null-trigger}
Under the unified minimality and protocol closure assumptions, mode axis completeness holds if and only if the Hankel rank is uniformly bounded:
there exist a constant \(d_\ast\) and a level threshold \(\ell_0\) such that for all \(\ell\ge \ell_0\) and all phase samples \(\theta\),
\[
d^{(\ell)}_{\min}(\theta)\le d_\ast.
\]
When the Hankel rank cannot be uniformly bounded, any fixed finite-dimensional mode type will exhibit non-cancellable modal collapse at some \((\ell,\theta)\); under protocol closure, this collapse cannot be silently absorbed and can only trigger recycling as a dimensional \(\mathsf{NULL}\).
\end{theorem}

\begin{proof}
This is the discussion-level restatement of Theorem \ref{thm:xi-completeness-iff-hankel-bounded}; its proof relies on Hankel rank Theorem \ref{thm:hankel-rank-minimal}. This completes the proof.
\end{proof}

\paragraph{Concluding synthesis (strict classification)}
Combining the commutative diagrams/indices/transversal uniqueness from this and the previous subsection, the mathematical substance of the ``fixed slice image recycler'' can be stated as: under the auditable caliber allowed by the regime, any difference can ultimately only manifest as
\begin{enumerate}
  \item the integer index \(\kappa=\deg(B_F)\) of pointwise defects (or the equivalent Toeplitz/Pontryagin/Tate--Maslov integer barrier);
  \item finite discrete residuals at the endpoint (e.g., \(m\bmod 4\));
  \item or, when continuous externalization is permitted, the unique radial transversal \(\varrho=\log|u|\) (and its monotone reparametrizations);
  \item together with finite field/congruence fingerprints compatible with the above compressions (Dwork guardrails and stable residual spectra, Corollary \ref{cor:discussion-recycling-residuals}).
\end{enumerate}
Beyond these residuals, all other degrees of freedom are at the type level either recycled as \(\mathsf{NULL}\) or compressed into verifiable primitive/exponential invariants.

\endinput
